\documentclass[numbers=enddot,12pt,final,onecolumn,notitlepage]{scrartcl}%
\usepackage[headsepline,footsepline,manualmark]{scrlayer-scrpage}
\usepackage[all,cmtip]{xy}
\usepackage{amssymb}
\usepackage{amsmath}
\usepackage{amsthm}
\usepackage{framed}
\usepackage{comment}
\usepackage{color}
\usepackage{hyperref}
\usepackage[sc]{mathpazo}
\usepackage[T1]{fontenc}
\usepackage{tikz}
\usepackage{needspace}
\usepackage{tabls}
\usepackage{wasysym}
\usepackage{easytable}
\usepackage{pythonhighlight}
%TCIDATA{OutputFilter=latex2.dll}
%TCIDATA{Version=5.50.0.2960}
%TCIDATA{LastRevised=Monday, November 22, 2021 11:52:10}
%TCIDATA{SuppressPackageManagement}
%TCIDATA{<META NAME="GraphicsSave" CONTENT="32">}
%TCIDATA{<META NAME="SaveForMode" CONTENT="1">}
%TCIDATA{BibliographyScheme=Manual}
%TCIDATA{Language=American English}
%BeginMSIPreambleData
\providecommand{\U}[1]{\protect\rule{.1in}{.1in}}
%EndMSIPreambleData
\usetikzlibrary{arrows.meta}
\usetikzlibrary{chains}
\newcounter{exer}
\newcounter{exera}
\numberwithin{exer}{subsection}
\theoremstyle{definition}
\newtheorem{theo}{Theorem}[subsection]
\newenvironment{theorem}[1][]
{\begin{theo}[#1]\begin{leftbar}}
{\end{leftbar}\end{theo}}
\newtheorem{lem}[theo]{Lemma}
\newenvironment{lemma}[1][]
{\begin{lem}[#1]\begin{leftbar}}
{\end{leftbar}\end{lem}}
\newtheorem{prop}[theo]{Proposition}
\newenvironment{proposition}[1][]
{\begin{prop}[#1]\begin{leftbar}}
{\end{leftbar}\end{prop}}
\newtheorem{defi}[theo]{Definition}
\newenvironment{definition}[1][]
{\begin{defi}[#1]\begin{leftbar}}
{\end{leftbar}\end{defi}}
\newtheorem{remk}[theo]{Remark}
\newenvironment{remark}[1][]
{\begin{remk}[#1]\begin{leftbar}}
{\end{leftbar}\end{remk}}
\newtheorem{coro}[theo]{Corollary}
\newenvironment{corollary}[1][]
{\begin{coro}[#1]\begin{leftbar}}
{\end{leftbar}\end{coro}}
\newtheorem{conv}[theo]{Convention}
\newenvironment{convention}[1][]
{\begin{conv}[#1]\begin{leftbar}}
{\end{leftbar}\end{conv}}
\newtheorem{quest}[theo]{Question}
\newenvironment{question}[1][]
{\begin{quest}[#1]\begin{leftbar}}
{\end{leftbar}\end{quest}}
\newtheorem{warn}[theo]{Warning}
\newenvironment{warning}[1][]
{\begin{warn}[#1]\begin{leftbar}}
{\end{leftbar}\end{warn}}
\newtheorem{conj}[theo]{Conjecture}
\newenvironment{conjecture}[1][]
{\begin{conj}[#1]\begin{leftbar}}
{\end{leftbar}\end{conj}}
\newtheorem{exam}[theo]{Example}
\newenvironment{example}[1][]
{\begin{exam}[#1]\begin{leftbar}}
{\end{leftbar}\end{exam}}
\newtheorem{exmp}[exer]{Exercise}
\newenvironment{exercise}[1][]
{\begin{exmp}[#1]\begin{leftbar}}
{\end{leftbar}\end{exmp}}
\newenvironment{statement}{\begin{quote}}{\end{quote}}
\newenvironment{fineprint}{\medskip \begin{small}}{\end{small} \medskip}
\iffalse
\newenvironment{proof}[1][Proof]{\noindent\textbf{#1.} }{\ \rule{0.5em}{0.5em}}
\newenvironment{question}[1][Question]{\noindent\textbf{#1.} }{\ \rule{0.5em}{0.5em}}
\newenvironment{warning}[1][Warning]{\noindent\textbf{#1.} }{\ \rule{0.5em}{0.5em}}
\newenvironment{teachingnote}[1][Teaching note]{\noindent\textbf{#1.} }{\ \rule{0.5em}{0.5em}}
\fi
\let\sumnonlimits\sum
\let\prodnonlimits\prod
\let\cupnonlimits\bigcup
\let\capnonlimits\bigcap
\renewcommand{\sum}{\sumnonlimits\limits}
\renewcommand{\prod}{\prodnonlimits\limits}
\renewcommand{\bigcup}{\cupnonlimits\limits}
\renewcommand{\bigcap}{\capnonlimits\limits}
\setlength\tablinesep{3pt}
\setlength\arraylinesep{3pt}
\setlength\extrarulesep{3pt}
\voffset=0cm
\hoffset=-0.7cm
\setlength\textheight{22.5cm}
\setlength\textwidth{15.5cm}
\newcommand\arxiv[1]{\href{http://www.arxiv.org/abs/#1}{\texttt{arXiv:#1}}}
\newenvironment{verlong}{}{}
\newenvironment{vershort}{}{}
\newenvironment{noncompile}{}{}
\newenvironment{teachingnote}{}{}
\excludecomment{verlong}
\includecomment{vershort}
\excludecomment{noncompile}
\excludecomment{teachingnote}
\newcommand{\CC}{\mathbb{C}}
\newcommand{\RR}{\mathbb{R}}
\newcommand{\QQ}{\mathbb{Q}}
\newcommand{\NN}{\mathbb{N}}
\newcommand{\ZZ}{\mathbb{Z}}
\newcommand{\KK}{\mathbb{K}}
\newcommand{\id}{\operatorname{id}}
\newcommand{\lcm}{\operatorname{lcm}}
\newcommand{\rev}{\operatorname{rev}}
\newcommand{\powset}[2][]{\ifthenelse{\equal{#2}{}}{\mathcal{P}\left(#1\right)}{\mathcal{P}_{#1}\left(#2\right)}}
\newcommand{\set}[1]{\left\{ #1 \right\}}
\newcommand{\abs}[1]{\left| #1 \right|}
\newcommand{\tup}[1]{\left( #1 \right)}
\newcommand{\ive}[1]{\left[ #1 \right]}
\newcommand{\floor}[1]{\left\lfloor #1 \right\rfloor}
\newcommand{\lf}[2]{#1^{\underline{#2}}}
\newcommand{\underbrack}[2]{\underbrace{#1}_{\substack{#2}}}
\newcommand{\horrule}[1]{\rule{\linewidth}{#1}}
\newcommand{\are}{\ar@{-}}
\newcommand{\nnn}{\nonumber\\}
\newcommand{\sslash}{\mathbin{/\mkern-6mu/}}
\newcommand{\numboxed}[2]{\underbrace{\boxed{#1}}_{\text{box } #2}}
\newcommand{\ig}[2]{\includegraphics[scale=#1]{#2.png}}
\newcommand{\UNFINISHED}{\begin{center} \Huge{\textbf{Unfinished material begins here.}} \end{center} }
\iffalse
\NOEXPAND{\today}{\today}
\NOEXPAND{\sslash}{\sslash}
\NOEXPAND{\numboxed}[2]{\numboxed}
\NOEXPAND{\UNFINISHED}{\UNFINISHED}
\fi
\ihead{Math 504 notes}
\ohead{page \thepage}
\cfoot{\today}
\begin{document}

\title{Math 504: Advanced Linear Algebra}
\author{Hugo Woerdeman, with edits by Darij Grinberg\thanks{Drexel University, Korman
Center, 15 S 33rd Street, Philadelphia PA, 19104, USA}}
\date{\today\ (unfinished!)}
\maketitle
\tableofcontents

\section*{Math 504 Lecture 22}

\section{Singular value decomposition}

This will be just a brief introduction to singular value decomposition. For
much more, see [Trefethen \& Bau].

\subsection{The $\operatorname*{Ker}\left(  A^{\ast}A\right)  $ lemma}

\begin{proposition}
[$\operatorname*{Ker}\left(  A^{\ast}A\right)  $ lemma]Let $A\in
\mathbb{C}^{m\times n}$ be any $m\times n$-matrix with complex entries (not
necessarily square).

Then:

\textbf{(a)} We have $\operatorname*{Ker}A=\operatorname*{Ker}\left(  A^{\ast
}A\right)  $.

\textbf{(b)} We have $\operatorname*{rank}A=\operatorname*{rank}\left(
A^{\ast}A\right)  $.
\end{proposition}

\begin{proof}
\textbf{(a)} Let $x\in\mathbb{C}^{n}$. Then, we have the following chain of
equivalences:%
\begin{align*}
\left(  x\in\operatorname*{Ker}A\right)  \   & \Longleftrightarrow\ \left(
Ax=0\right)  \\
& \Longleftrightarrow\ \left(  \left\vert \left\vert Ax\right\vert \right\vert
=0\right)  \ \ \ \ \ \ \ \ \ \ \left(  \text{since a vector is }0\text{ if and
only if its length is }0\right)  \\
& \Longleftrightarrow\ \left(  \sqrt{\left(  Ax\right)  ^{\ast}Ax}=0\right)
\ \ \ \ \ \ \ \ \ \ \left(  \text{since }\left\vert \left\vert Ax\right\vert
\right\vert =\sqrt{\left\langle Ax,\ Ax\right\rangle }=\sqrt{\left(
Ax\right)  ^{\ast}Ax}\right)  \\
& \Longleftrightarrow\ \left(  \left(  Ax\right)  ^{\ast}Ax=0\right)  \\
& \Longleftrightarrow\ \left(  x^{\ast}A^{\ast}Ax=0\right)  \\
& \Longleftarrow\ \left(  A^{\ast}Ax=0\right)  \\
& \Longleftrightarrow\ \left(  x\in\operatorname*{Ker}\left(  A^{\ast
}A\right)  \right)  .
\end{align*}
So we have shown that $\left(  x\in\operatorname*{Ker}\left(  A^{\ast
}A\right)  \right)  \ \Longrightarrow\ \left(  x\in\operatorname*{Ker}%
A\right)  $. But the converse implication is even easier:%
\[
\left(  x\in\operatorname*{Ker}A\right)  \ \Longleftrightarrow\ \left(
Ax=0\right)  \ \Longrightarrow\ \left(  A^{\ast}Ax=0\right)
\ \Longleftrightarrow\ \left(  x\in\operatorname*{Ker}\left(  A^{\ast
}A\right)  \right)  .
\]
So we get the equivalence%
\[
\left(  x\in\operatorname*{Ker}A\right)  \ \Longleftrightarrow\ \left(
x\in\operatorname*{Ker}\left(  A^{\ast}A\right)  \right)  .
\]
Since this holds for all vectors $x\in\mathbb{C}^{n}$, we thus conclude that
$\operatorname*{Ker}A=\operatorname*{Ker}\left(  A^{\ast}A\right)  $. This
proves part \textbf{(a)}.

\textbf{(b)} The rank-nullity theorem yields%
\begin{align*}
\operatorname*{rank}A  & =n-\dim\left(  \operatorname*{Ker}A\right)
\ \ \ \ \ \ \ \ \ \ \text{and}\\
\operatorname*{rank}\left(  A^{\ast}A\right)    & =n-\dim\left(
\operatorname*{Ker}\left(  A^{\ast}A\right)  \right)  .
\end{align*}
The right hand sides of these two equalities are equal (since part
\textbf{(a)} yields $\operatorname*{Ker}A=\operatorname*{Ker}\left(  A^{\ast
}A\right)  $). Thus, the left hand sides are also equal. In other words,
$\operatorname*{rank}A=\operatorname*{rank}\left(  A^{\ast}A\right)  $. This
proves part \textbf{(b)}.
\end{proof}

Note that the above proposition really requires a matrix with complex entries.
It cannot be generalized to matrices over an arbitrary field.

\subsection{The singular value decomposition}

\begin{definition}
Let $A$ and $B$ be two $m\times n$-matrices with complex entries.

We say that $A$ and $B$ are \textbf{unitarily equivalent} if there exist
unitary matrices $U\in\operatorname*{U}\nolimits_{m}\left(  \mathbb{C}\right)
$ and $V\in\operatorname*{U}\nolimits_{n}\left(  \mathbb{C}\right)  $ such
that $A=UBV^{\ast}$.
\end{definition}

Note the difference between \textquotedblleft unitarily
equivalent\textquotedblright\ and \textquotedblleft unitarily
similar\textquotedblright: The latter requires $A=UBU^{\ast}$, whereas the
former only requires $A=UBV^{\ast}$.

Unitary equivalence is an equivalence relation (exercise!). A natural question
is therefore: What is the \textquotedblleft simplest\textquotedblright\ matrix
in the equivalence class of a matrix?

\begin{definition}
A rectangular matrix $A\in\mathbb{F}^{n\times m}$ (for a field $\mathbb{F}$)
is said to be \textbf{pseudodiagonal} if it satisfies%
\[
A_{i,j}=0\ \ \ \ \ \ \ \ \ \ \text{whenever }i\neq j.
\]

\end{definition}

This is just the straightforward generalization of diagonal matrices to
non-square matrices. In particular, a square matrix is pseudodiagonal if and
only if it is diagonal. A pseudodiagonal $2\times3$-matrix looks like this:
$\left(
\begin{array}
[c]{ccc}%
\ast & 0 & 0\\
0 & \ast & 0
\end{array}
\right)  $. A pseudodiagonal $3\times2$-matrix looks like this: $\left(
\begin{array}
[c]{cc}%
\ast & 0\\
0 & \ast\\
0 & 0
\end{array}
\right)  $. (Of course, any of the $\ast$s can be a $0$ too.)

\begin{theorem}
[SVD]Let $A\in\mathbb{C}^{m\times n}$. Then:

\textbf{(a)} There exist unitary matrices $U\in\operatorname*{U}%
\nolimits_{m}\left(  \mathbb{C}\right)  $ and $V\in\operatorname*{U}%
\nolimits_{n}\left(  \mathbb{C}\right)  $ and a pseudodiagonal matrix $\Sigma$
such that%
\[
A=U\Sigma V^{\ast}.
\]
In other words, $A$ is unitarily equivalent to a pseudodiagonal matrix.

\textbf{(b)} The matrix $\Sigma$ is unique up to permutation of its diagonal
entries. (The matrices $U$ and $V$ are usually not unique.)

\textbf{(c)} Let $k=\operatorname*{rank}A$. Then, the matrix $\Sigma$ has
exactly $k$ nonzero diagonal entries.

\textbf{(d)} Let $\sigma_{1},\sigma_{2},\ldots,\sigma_{n}$ be the square roots
of the eigenvalues of the Hermitian matrix $A^{\ast}A$, listed in decreasing
order (so that $\sigma_{1}\geq\sigma_{2}\geq\cdots\geq\sigma_{m}$). Then, we
have $\sigma_{k+1}=\sigma_{k+2}=\cdots=\sigma_{m}=0$, and we can take
\[
\Sigma=\left(
\begin{array}
[c]{ccccccc}%
\sigma_{1} & 0 & \cdots & 0 & 0 & \cdots & 0\\
0 & \sigma_{2} & \cdots & 0 & 0 & \cdots & 0\\
\vdots & \vdots & \ddots & \vdots & \vdots & \ddots & \vdots\\
0 & 0 & \cdots & \sigma_{k} & 0 & \cdots & 0\\
0 & 0 & \cdots & 0 & 0 & \cdots & 0\\
\vdots & \vdots & \ddots & \vdots & \vdots & \ddots & \vdots\\
0 & 0 & \cdots & 0 & 0 & \cdots & 0
\end{array}
\right)  \in\mathbb{C}^{m\times n}%
\]
in part \textbf{(a)}.
\end{theorem}

\begin{definition}
The triple $\left(  U,V,\Sigma\right)  $ in part \textbf{(a)} of the theorem
is called a \textbf{singular value decomposition} (short: \textbf{SVD}) of
$A$. The numbers $\sigma_{1},\sigma_{2},\ldots,\sigma_{n}$ are called the
\textbf{singular values} of $A$.
\end{definition}

Before we prove the theorem, a few words about the use of an SVD. Often in
practice, you have a matrix $A$ and you want to find a low-rank
\textquotedblleft approximation\textquotedblright\ for $A$: that is, a matrix
$B$ that is \textquotedblleft sufficiently close\textquotedblright\ to $A$ and
yet has low rank. Indeed, the idea for doing this is to compute an SVD of $A$
-- that is, $A=U\Sigma V^{\ast}$ with $U,\Sigma,V$ as above -- and then set
all but the first few $\sigma_{i}$'s to $0$ in $\Sigma$.

Let us now prove the theorem.

\begin{proof}
[Proof of the SVD theorem.] Let $k=\operatorname*{rank}A$. Then,
$k=\operatorname*{rank}\left(  A^{\ast}A\right)  $ (as we know from part
\textbf{(b)} of the above proposition). Also, $k$ is the rank of an $m\times
n$-matrix (namely, $A$), and thus satisfies $k\leq m$ and $k\leq n$.

The matrix $A^{\ast}A$ is Hermitian, since $\left(  A^{\ast}A\right)  ^{\ast
}=A^{\ast}\underbrace{\left(  A^{\ast}\right)  ^{\ast}}_{=A}=A^{\ast}A$.
Moreover, this matrix $A^{\ast}A$ is positive semidefinite, since each vector
$x\in\mathbb{C}^{n}$ satisfies%
\[
\left\langle A^{\ast}Ax,\ x\right\rangle =\underbrace{x^{\ast}A^{\ast}%
}_{=\left(  Ax\right)  ^{\ast}}Ax=\left(  Ax\right)  ^{\ast}Ax=\left\vert
\left\vert Ax\right\vert \right\vert ^{2}\geq0.
\]
The spectral theorem thus yields that $A^{\ast}A$ can be written in the form
\[
A^{\ast}A=V\cdot\operatorname*{diag}\left(  \lambda_{1},\lambda_{2}%
,\ldots,\lambda_{n}\right)  \cdot V^{\ast},
\]
where $V\in\operatorname*{U}\nolimits_{n}\left(  \mathbb{C}\right)  $ is a
unitary matrix and where the $\lambda_{1},\lambda_{2},\ldots,\lambda_{n}$ are
the eigenvalues of $A^{\ast}A$ listed in decreasing order. Consider this $V$
and these $\lambda_{1},\lambda_{2},\ldots,\lambda_{n}$. The eigenvalues
$\lambda_{1},\lambda_{2},\ldots,\lambda_{n}$ are nonnegative reals, since
$A^{\ast}A$ is positive semidefinite. Thus, we can set%
\[
\sigma_{i}:=\sqrt{\lambda_{i}}\ \ \ \ \ \ \ \ \ \ \text{for each }i\in\left[
n\right]  .
\]
Thus, $\sigma_{1},\sigma_{2},\ldots,\sigma_{n}$ are nonnegative reals.

We have%
\begin{align*}
k  & =\operatorname*{rank}\left(  A^{\ast}A\right)  =\operatorname*{rank}%
\left(  \operatorname*{diag}\left(  \lambda_{1},\lambda_{2},\ldots,\lambda
_{n}\right)  \right)  \ \ \ \ \ \ \ \ \ \ \left(  \text{since }A^{\ast
}A\overset{\operatorname*{us}}{\sim}\operatorname*{diag}\left(  \lambda
_{1},\lambda_{2},\ldots,\lambda_{n}\right)  \right)  \\
& =\left(  \text{the number of }i\in\left[  n\right]  \text{ such that
}\lambda_{i}\neq0\right)  \\
& =\left(  \text{the number of }i\in\left[  n\right]  \text{ such that }%
\sigma_{i}\neq0\right)  \ \ \ \ \ \ \ \ \ \ \left(  \text{since }\sigma
_{i}=\sqrt{\lambda_{i}}\right)  .
\end{align*}
Hence, exactly $k$ of the numbers $\sigma_{1},\sigma_{2},\ldots,\sigma_{n}$
are nonzero. Since $\sigma_{1},\sigma_{2},\ldots,\sigma_{n}$ are nonnegative
reals and sorted in weakly decreasing order, this entails that%
\begin{align*}
\sigma_{1}  & \geq\sigma_{2}\geq\cdots\geq\sigma_{k}%
>0\ \ \ \ \ \ \ \ \ \ \text{and}\\
\sigma_{k+1}  & =\sigma_{k+2}=\cdots=\sigma_{n}=0.
\end{align*}


Let $v_{1},v_{2},\ldots,v_{n}$ be the columns of the unitary matrix $V$. Thus,
$\left(  v_{1},v_{2},\ldots,v_{n}\right)  $ is an orthonormal basis of
$\mathbb{C}^{n}$ (since $V$ is unitary). Since $V$ is the unitary matrix in a
spectral decomposition of $A^{\ast}A$, the columns of $V$ are eigenvectors of
$A^{\ast}A$, corresponding to the eigenvalues $\lambda_{1},\lambda_{2}%
,\ldots,\lambda_{n}$. In other words,%
\[
A^{\ast}Av_{i}=\lambda_{i}v_{i}\ \ \ \ \ \ \ \ \ \ \text{for each }i\in\left[
n\right]  .
\]


For each $j\in\left[  k\right]  $, set%
\[
u_{j}:=\dfrac{1}{\sigma_{j}}Av_{j}.
\]
(The division by $\sigma_{j}$ is well-defined, since $\sigma_{1}\geq\sigma
_{2}\geq\cdots\geq\sigma_{k}>0$ imply $\sigma_{j}\neq0$.) We claim that the
tuple $\left(  u_{1},u_{2},\ldots,u_{k}\right)  $ is orthonormal. Indeed:

\begin{itemize}
\item For any $i\neq j$, we have%
\begin{align*}
\left\langle u_{i},u_{j}\right\rangle  & =\left\langle \dfrac{1}{\sigma_{i}%
}Av_{i},\ \dfrac{1}{\sigma_{j}}Av_{j}\right\rangle =\dfrac{1}{\sigma_{i}%
\sigma_{j}}\left\langle Av_{i},\ Av_{j}\right\rangle =\dfrac{1}{\sigma
_{i}\sigma_{j}}\left(  Av_{j}\right)  ^{\ast}Av_{i}\\
& =\dfrac{1}{\sigma_{i}\sigma_{j}}v_{j}^{\ast}\underbrace{A^{\ast}Av_{i}%
}_{=\lambda_{i}v_{i}}=\dfrac{1}{\sigma_{i}\sigma_{j}}v_{j}^{\ast}\lambda
_{i}v_{i}=\dfrac{\lambda_{i}}{\sigma_{i}\sigma_{j}}\underbrace{v_{j}^{\ast
}v_{i}}_{\substack{=\left\langle v_{i},v_{j}\right\rangle \\=0\\\text{(since
}\left(  v_{1},v_{2},\ldots,v_{n}\right)  \\\text{is orthonormal)}}}=0,
\end{align*}
and thus $u_{i}\perp u_{j}$.

\item For any $i$, we have%
\begin{align*}
\left\langle u_{i},u_{i}\right\rangle  & =\left\langle \dfrac{1}{\sigma_{i}%
}Av_{i},\ \dfrac{1}{\sigma_{i}}Av_{i}\right\rangle =\dfrac{1}{\sigma_{i}%
\sigma_{i}}\left\langle Av_{i},\ Av_{i}\right\rangle =\dfrac{1}{\sigma
_{i}\sigma_{i}}\left(  Av_{i}\right)  ^{\ast}Av_{i}\\
& =\dfrac{1}{\sigma_{i}\sigma_{i}}v_{i}^{\ast}\underbrace{A^{\ast}Av_{i}%
}_{=\lambda_{i}v_{i}}=\underbrace{\dfrac{\lambda_{i}}{\sigma_{i}\sigma_{i}}%
}_{\substack{=1\\\text{(since }\sigma_{i}=\sqrt{\lambda_{i}}\text{)}}%
}v_{i}^{\ast}v_{i}=v_{i}^{\ast}v_{i}=\left\langle v_{i},v_{i}\right\rangle =1
\end{align*}
(since $\left(  v_{1},v_{2},\ldots,v_{n}\right)  $ is orthonormal), and thus
$\left\vert \left\vert u_{i}\right\vert \right\vert =1$.
\end{itemize}

So $\left(  u_{1},u_{2},\ldots,u_{k}\right)  $ is an orthonormal tuple of
vectors in $\mathbb{C}^{m}$. As we know, we can extend it to an orthonormal
basis $\left(  u_{1},u_{2},\ldots,u_{m}\right)  $ of $\mathbb{C}^{m}$. So let
us do this. Let $U$ be the $m\times m$-matrix with columns $u_{1},u_{2}%
,\ldots,u_{m}$. Thus, $U$ is a unitary matrix, since its columns form an
orthonormal basis of $\mathbb{C}^{m}$.

Set $\sigma_{i}:=0$ for every $i>n$. Note that $\sigma_{i}=0$ for every $i>k$,
since we already have $\sigma_{k+1}=\sigma_{k+2}=\cdots=\sigma_{n}=0$.

Let%
\[
\Sigma:=\left(
\begin{array}
[c]{ccccccc}%
\sigma_{1} & 0 & \cdots & 0 & 0 & \cdots & 0\\
0 & \sigma_{2} & \cdots & 0 & 0 & \cdots & 0\\
\vdots & \vdots & \ddots & \vdots & \vdots & \ddots & \vdots\\
0 & 0 & \cdots & \sigma_{k} & 0 & \cdots & 0\\
0 & 0 & \cdots & 0 & 0 & \cdots & 0\\
\vdots & \vdots & \ddots & \vdots & \vdots & \ddots & \vdots\\
0 & 0 & \cdots & 0 & 0 & \cdots & 0
\end{array}
\right)  \in\mathbb{C}^{m\times n}.
\]


Now, we claim that $A=U\Sigma V^{\ast}$. To prove this, we shall show the
equivalent claim $AV=U\Sigma$.

It is sufficient to prove that
\[
\left(  AV\right)  _{\bullet,j}=\left(  U\Sigma\right)  _{\bullet
,j}\ \ \ \ \ \ \ \ \ \ \text{for each }j\in\left[  n\right]
\]
(recall that the notation $M_{\bullet,j}$ means the $j$-th column of a matrix
$M$). So let us fix $j\in\left[  n\right]  $ and try to prove that $\left(
AV\right)  _{\bullet,j}=\left(  U\Sigma\right)  _{\bullet,j}$.

\begin{itemize}
\item If $j\leq k$, then we have%
\begin{align*}
\left(  AV\right)  _{\bullet,j}  & =A\underbrace{V_{\bullet,j}}_{=v_{j}%
}=Av_{j}=\sigma_{j}\underbrace{u_{j}}_{=U_{\bullet,j}}%
\ \ \ \ \ \ \ \ \ \ \left(  \text{since }u_{j}=\dfrac{1}{\sigma_{j}}%
Av_{j}\right)  \\
& =\sigma_{j}U_{\bullet,j}=\left(  U\Sigma\right)  _{\bullet,j}.
\end{align*}


\item If $j>k$, then we can apply the $A^{\ast}Av_{i}=\lambda_{i}v_{i}$
equality to $i=j$ and obtain%
\[
A^{\ast}Av_{j}=\underbrace{\lambda_{j}}_{\substack{=0\\\text{(since
}j>k\text{)}}}v_{j}=0,
\]
so that $v_{j}\in\operatorname*{Ker}\left(  A^{\ast}A\right)
=\operatorname*{Ker}A$ (by the proposition above); now,%
\[
\left(  AV\right)  _{\bullet,j}=A\underbrace{V_{\bullet,j}}_{=v_{j}}%
=Av_{j}=0\ \ \ \ \ \ \ \ \ \ \left(  \text{since }v_{j}=\operatorname*{Ker}%
A\right)  ,
\]
but also%
\[
\left(  U\Sigma\right)  _{\bullet,j}=\underbrace{\sigma_{j}}_{=\sqrt
{\lambda_{j}}=\sqrt{0}=0}U_{\bullet,j}=0,
\]
and therefore $\left(  AV\right)  _{\bullet,j}=0=\left(  U\Sigma\right)
_{\bullet,j}$.
\end{itemize}

Thus, we have proved $\left(  AV\right)  _{\bullet,j}=\left(  U\Sigma\right)
_{\bullet,j}$ in both cases, and this completes the proof of $AV=U\Sigma$.
Therefore,%
\[
\underbrace{U\Sigma}_{=AV}V^{\ast}=A\underbrace{VV^{\ast}}_{\substack{=I_{n}%
\\\text{(since }V\text{ is unitary)}}}=A,
\]
so that $A=U\Sigma V^{\ast}$, as desired.

This proves parts \textbf{(a)}, \textbf{(c)} and \textbf{(d)} of the theorem.

\textbf{(b)} We must show that if $P$ is a pseudodiagonal matrix such that $A$
is unitarily equivalent to $P$, then $P$ and $\Sigma$ have the same diagonal
entries up to order.

Rough outline: The singular values of a matrix are invariant under unitary
equivalence. In other words, if $A=UBV^{\ast}$ for some unitary $U$ and $V$,
then $A$ and $B$ have the same singular values. (In fact, $A=UBV^{\ast}$
entails%
\[
A^{\ast}A=\left(  UBV^{\ast}\right)  ^{\ast}UBV^{\ast}=VB^{\ast}%
\underbrace{U^{\ast}U}_{=I_{m}}BV^{\ast}=VB^{\ast}BV^{\ast}%
\overset{\operatorname*{us}}{\sim}B^{\ast}B;
\]
this shows that $A^{\ast}A$ and $B^{\ast}B$ have the same eigenvalues;
therefore $A$ and $B$ have the same singular values.) Thus, $P$ and $\Sigma$
must have the same singular values (since $P$ and $\Sigma$ are both unitarily
equivalent to $A$). However, for a pseudodiagonal matrix, the singular values
are simply the diagonal entries. (Fine print: There might be $0$s not
accounted for.)
\end{proof}


\end{document}